Considera la funció
\[
f(x)=x^3+2x^2-1 .
\]

\begin{apartats}

\apartat[1,25]{1}
Enuncia el teorema de Bolzano i demostra que l'equació $f(x)=0$ té almenys una solució a
l'interval $[0,1]$.

\begin{solucio}
\textbf{Teorema de Bolzano.} Si $f$ és contínua a $[a,b]$ i $f(a)$ i $f(b)$ tenen signes
diferents, aleshores existeix almenys un $c\in(a,b)$ amb $f(c)=0$.\\
$f$ és polinòmica i, per tant, contínua a $[0,1]$. A més, $f(0)=-1<0$ i $f(1)=1+2-1=2>0$.
Com que canvia de signe, hi ha almenys una solució a $(0,1)$.
\end{solucio}

\apartat[1,25]{0,75}
Aplicant el mètode de la bisecció, troba un interval de longitud $0{,}25$ que contingui una
solució de l'equació. Justifica cada pas.

\begin{solucio}
$f(0{,}5)=0{,}125+0{,}5-1=-0{,}375<0$: com que $f(1)>0$, l'arrel és a
$\left(0{,}5;\,1\right)$.\\
$f(0{,}75)=0{,}421875+1{,}125-1=0{,}546875>0$: com que $f(0{,}5)<0$, l'arrel és a
$\left(0{,}5;\,0{,}75\right)$, de longitud $0{,}25$.
\end{solucio}

\begin{nomesllarg}
\apartat{0,75}
Demostra que les gràfiques de les funcions $y=\ln x$ i $y=2-x$ es tallen en algun punt
d'abscissa $x\in(1,2)$.

\begin{solucio}
Es tallen on $\ln x=2-x$, és a dir on s'anul·la $k(x)=\ln x+x-2$, que és contínua a $(0,+\infty)$
i, per tant, a $[1,2]$.\\
$k(1)=0+1-2=-1<0$ i $k(2)=\ln2+0\approx0{,}69>0$. Per Bolzano, hi ha un punt de tall a $(1,2)$.
\end{solucio}
\end{nomesllarg}

\end{apartats}
